\section{Recognizing the correct counting model}

The following diagnostic table summarizes the main finite counting structures.

\begin{center}
\begin{tabular}{llll}
\toprule
Outcome type & Order? & Repetition? & Count \\
\midrule
sequence with repetition & yes & yes & \(n^k\) \\
ordered selection & yes & no & \(P(n,k)=\dfrac{n!}{(n-k)!}\) \\
permutation & yes & no, all used & \(n!\) \\
circular arrangement & relative order & no, all used & \((n-1)!\) \\
combination & no & no & \(\binom{n}{k}\) \\
combination with repetition & no & yes & \(\binom{n+k-1}{k}\) \\
multiset permutation & positions yes & identical objects & \(\dfrac{n!}{n_1!\cdots n_r!}\) \\
\bottomrule
\end{tabular}
\end{center}

The table should be used only after the outcome type has been identified. It is a summary, not a substitute for modelling.

\begin{remarkbox}
A formula is a compressed argument. If the argument does not match the situation, the formula gives the wrong count even when the arithmetic is correct.
\end{remarkbox}

\section{Common traps}

\subsection*{Trap 1: confusing a sequence with a set}

The ordered outcomes
\[
(A,B,C)\quad\text{and}\quad(C,B,A)
\]
are different sequences but the same set:
\[
\{A,B,C\}=\{C,B,A\}.
\]
Use sequences when order is recorded. Use sets when only membership is recorded.

\subsection*{Trap 2: treating indistinguishable objects as distinguishable}

The letters in \(\text{BANANA}\) are not six distinct letters. The three \(A\)'s are indistinguishable, and the two \(N\)'s are indistinguishable. Therefore \(6!\) overcounts the number of visible arrangements.

\subsection*{Trap 3: using combinations when roles matter}

A committee is a set. A medal assignment is an ordered assignment. A group of three athletes and a gold-silver-bronze result are not the same kind of outcome.

\subsection*{Trap 4: using a representation that records too much}

Sometimes an auxiliary representation is more detailed than the intended outcome. For example, ordered lists may be useful for deriving the number of committees, but the committee itself is not an ordered list. If the extra descriptions are not corrected for, the count is too large.

\subsection*{Trap 5: using a representation that records too little}

Sometimes a representation hides distinctions that the model requires. If positions, roles, labels, or order are part of the outcome, then a representation that ignores them gives too small a count. A good representation must preserve exactly the distinctions that the recording rule declares visible.
